\section*{Original bearing operations (pp.9:6--7)}
Given $A,B\subseteq\mathbb R^n$, their Minkowski sum is defined as
\[A+B:=\{a+b:a\in A,b\in B\}.\]
If $P$ and $Q$ are polytopes then $P+Q$ is also a polytope. In two-dimensions, Minkowski sum has nice properties. Let $P$ be a polytope in $\mathbb R^2$ with vertices $v_1,\ldots,v_k$ where $k>1$. We can assume they are ordered so that $P$ lies in the left closed half plane determined by the line going from $v_i$ to $v_{i+1}$ for $i<k$, and similarly for $v_k$ and $v_1$. Let $E(P)$ be the collection of unit vectors in the direction of these $k$ edges. That is, vectors of the form $(v_{i+1}-v_i)/\|v_{i+1}-v_i\|$ for $i<k$, and $(v_1-v_k)/\|v_1-v_k\|$. If $|\vertset(P)|\le1$ then $E(P):=\emptyset$.
\begin{sourceresult}{Lemma 8}
Let $P_1,\ldots,P_r$ be non-empty polytopes in $\mathbb R^2$. Then $E(P_1+\cdots+P_r)=\bigcup_{i=1}^r E(P_i)$. Consequently, $|\vertset(P_1+\cdots+P_r)|\le\sum_{i=1}^r|\vertset(P_i)|$. The latter holds for empty $P_i$'s as well.
\end{sourceresult}
The lemma is folklore. It can be inferred from Chapter 13.3 in [11], and we give only an outline of proof.
\begin{proof}[Proof sketch of Lemma 8]
Given a non-empty polytope $P$ and $w\in\mathbb R^2$, let $P^w:=\{x\in P:\langle x,w\rangle=\max_{z\in P}\langle z,w\rangle\}$ be the set of extreme of points of $P$ in the direction $w$. It is either a vertex or an edge of $P$. For a pair of polytopes we have $(P_1+P_2)^w=P_1^w+P_2^w$. Every edge of $P_1$ yields an edge of $P_1+P_2$ with the same direction. Conversely, every edge of $P_1+P_2$ comes from one of $P_1$ or $P_2$.
\end{proof}
The second operation we use is
\[A\sqcup B:=\conv(A\cup B).\]
If $P$ and $Q$ are polytopes then $P\sqcup Q$ is also a polytope.
\begin{sourceresult}{Lemma 11}
Let $f,g$ be Laurent polynomials over $\mathbb F$.
\begin{enumerate}
\item[(i)] Then $\Newt(fg)=\Newt(f)+\Newt(g)$.
\item[(ii)] $\Newt(f+g)=\Newt(f)\sqcup\Newt(g)$, provided $\mathbb F=\mathbb R$ and both $f$ and $g$ are monotone.
\end{enumerate}
\end{sourceresult}
\begin{proof}
Part (i) can be found in [16] for polynomials; it extends to Laurent polynomials. Part (ii) is straightforward to verify.
\end{proof}
